{\scriptsize
\setlength{\tabcolsep}{3pt}
\begin{longtable}{>{\raggedright\arraybackslash}p{1.10cm}>{\raggedright\arraybackslash}p{1.80cm}>{\raggedright\arraybackslash}p{1.30cm}>{\raggedright\arraybackslash}p{1.60cm}>{\raggedright\arraybackslash}p{2.30cm}>{\raggedright\arraybackslash}p{2.30cm}>{\raggedright\arraybackslash}p{1.70cm}>{\raggedright\arraybackslash}p{1.70cm}>{\raggedright\arraybackslash}p{1.00cm}}
\caption{Cas climatique, sous-échantillons espacés ($D=5$, référence Köppen--Geiger seule) : moyennes sur les réplications}\label{tab:climat_sous}\\
\toprule
\textbf{$N$} & \textbf{Bras} & \textbf{$\hat K$ moy.} & \textbf{ARI Köppen} & \textbf{LL hors bloc (méd.)} & \textbf{LL prédictive VB (méd.)} & \textbf{Dégénérés} & \textbf{Temps méd. (s)} & \textbf{Rép.} \\
\midrule
\endfirsthead
\toprule
\textbf{$N$} & \textbf{Bras} & \textbf{$\hat K$ moy.} & \textbf{ARI Köppen} & \textbf{LL hors bloc (méd.)} & \textbf{LL prédictive VB (méd.)} & \textbf{Dégénérés} & \textbf{Temps méd. (s)} & \textbf{Rép.} \\
\midrule
\endhead
\bottomrule
\endfoot
30 & EM-BIC & 9,5 & 0,129 & -2123,72 & n/d & 100\,\% & 1,59 & 40 \\
30 & EM-ICL & 9,5 & 0,129 & -2123,72 & n/d & 100\,\% & 1,61 & 40 \\
30 & EM-K & 5,0 & 0,134 & -146,21 & n/d & 100\,\% & 0,19 & 40 \\
30 & EM-Kmax & 10,0 & 0,120 & -2065,31 & n/d & 100\,\% & 0,19 & 40 \\
30 & MAP-EM & 9,8 & 0,115 & -14,26 & n/d & 100\,\% & 1,95 & 40 \\
30 & VB-diff & 5,1 & 0,166 & -4,60 & -4,13 & 100\,\% & 1,39 & 40 \\
30 & VB-emp & 6,7 & 0,151 & -7,57 & -3,46 & 100\,\% & 1,40 & 40 \\
60 & EM-BIC & 9,7 & 0,132 & -109,78 & n/d & 100\,\% & 2,80 & 40 \\
60 & EM-ICL & 9,7 & 0,132 & -109,78 & n/d & 100\,\% & 2,98 & 40 \\
60 & EM-K & 5,0 & 0,154 & -18,57 & n/d & 100\,\% & 0,32 & 40 \\
60 & EM-Kmax & 10,0 & 0,133 & -104,84 & n/d & 100\,\% & 0,33 & 40 \\
60 & MAP-EM & 9,9 & 0,123 & -7,95 & n/d & 100\,\% & 1,68 & 40 \\
60 & VB-diff & 5,5 & 0,159 & -3,75 & -3,46 & 100\,\% & 1,62 & 40 \\
60 & VB-emp & 6,8 & 0,145 & -4,94 & -2,93 & 98\,\% & 1,51 & 40 \\
120 & EM-BIC & 7,3 & 0,155 & -19,74 & n/d & 90\,\% & 6,82 & 40 \\
120 & EM-ICL & 7,4 & 0,160 & -19,08 & n/d & 90\,\% & 6,79 & 40 \\
120 & EM-K & 5,0 & 0,145 & -7,98 & n/d & 62\,\% & 0,75 & 40 \\
120 & EM-Kmax & 10,0 & 0,141 & -27,55 & n/d & 100\,\% & 1,02 & 40 \\
120 & MAP-EM & 9,8 & 0,107 & -6,53 & n/d & 95\,\% & 2,76 & 40 \\
120 & VB-diff & 5,2 & 0,147 & -3,50 & -3,04 & 70\,\% & 2,89 & 40 \\
120 & VB-emp & 6,7 & 0,126 & -4,12 & -2,61 & 98\,\% & 2,49 & 40 \\
500 & EM-BIC & 9,6 & 0,155 & -0,33 & n/d & 22\,\% & 21,22 & 40 \\
500 & EM-ICL & 9,5 & 0,156 & -0,33 & n/d & 20\,\% & 21,07 & 40 \\
500 & EM-K & 5,0 & 0,127 & -1,08 & n/d & 0\,\% & 1,76 & 40 \\
500 & EM-Kmax & 10,0 & 0,150 & -0,52 & n/d & 20\,\% & 3,98 & 40 \\
500 & MAP-EM & 10,0 & 0,096 & -1,11 & n/d & 10\,\% & 6,37 & 40 \\
500 & VB-diff & 7,2 & 0,147 & -1,82 & -1,79 & 50\,\% & 4,80 & 40 \\
500 & VB-emp & 8,8 & 0,106 & -0,95 & -0,95 & 57\,\% & 5,21 & 40 \\
2000 & EM-BIC & 10,0 & 0,140 & 0,49 & n/d & 0\,\% & 38,11 & 40 \\
2000 & EM-ICL & 10,0 & 0,140 & 0,49 & n/d & 0\,\% & 38,44 & 40 \\
2000 & EM-K & 5,0 & 0,129 & -0,54 & n/d & 0\,\% & 2,60 & 40 \\
2000 & EM-Kmax & 10,0 & 0,140 & 0,49 & n/d & 0\,\% & 7,95 & 40 \\
2000 & MAP-EM & 10,0 & 0,106 & 0,15 & n/d & 0\,\% & 10,78 & 40 \\
2000 & VB-diff & 9,6 & 0,136 & -0,04 & -0,05 & 20\,\% & 8,96 & 40 \\
2000 & VB-emp & 10,0 & 0,106 & 0,14 & 0,13 & 0\,\% & 10,08 & 40 \\
complet & EM-BIC & 10,0 & 0,145 & 0,61 & n/d & 0\,\% & 693,07 & 1 \\
complet & EM-ICL & 10,0 & 0,145 & 0,61 & n/d & 0\,\% & 477,83 & 1 \\
complet & EM-K & 5,0 & 0,094 & -0,50 & n/d & 0\,\% & 27,25 & 1 \\
complet & EM-Kmax & 10,0 & 0,145 & 0,61 & n/d & 0\,\% & 122,31 & 1 \\
complet & MAP-EM & 10,0 & 0,141 & 0,54 & n/d & 0\,\% & 108,67 & 1 \\
complet & VB-diff & 10,0 & 0,145 & 0,52 & 0,52 & 0\,\% & 103,40 & 1 \\
complet & VB-emp & 10,0 & 0,141 & 0,54 & 0,54 & 0\,\% & 89,81 & 1 \\
\end{longtable}
\par\noindent{\footnotesize LL : log-vraisemblance moyenne par observation sur des blocs de 100 km sans point d'apprentissage ; $\hat K$ des bras MAP-EM et VB : composantes avec $N_k/N\geq1\,\%$. Une seule exécution à $N$ complet.}
}
